\section{Proofs}\label{app:proofs}
We identify $g=(G_1,\dots,G_L)$ with a vector in $\R^d$ and use the Euclidean inner product.

\subsection{Proof of Theorem~\ref{thm:nogo}}
Condition on the probes, which are independent of $g$. Let $V=\mathrm{span}\{E_1,\dots,E_N\}$, $k=\dim V\le N$, and let $\Pi$ be the orthogonal projector on $V$. The data $y_j=\ip{E_j}{g}=\ip{E_j}{\Pi g}$ are a function of $\Pi g$. The map $Q=2\Pi-I$ is orthogonal, so $Qg\overset d=g$, and $\Pi Qg=\Pi g$, $(I-\Pi)Qg=-(I-\Pi)g$. Hence
\[
\big(\Pi g,\,(I-\Pi)g\big)\overset d=\big(\Pi g,\,-(I-\Pi)g\big),\qquad\text{so}\qquad\E\big[(I-\Pi)g\mid\Pi g\big]=0 .
\]
Let $\tilde g$ be any estimator; given the probes it is a (possibly randomized) function of $\Pi g$. Write $\tilde g-g=(\tilde g-\Pi g)-(I-\Pi)g$. The cross term has mean zero by the last display, so
\[
\E\|\tilde g-g\|^2=\E\|\tilde g-\Pi g\|^2+\E\|(I-\Pi)g\|^2\ \ge\ \E\|(I-\Pi)g\|^2=\Big(1-\frac kd\Big)\E\|g\|^2\ \ge\ 1-\frac Nd .
\]
The last equality uses $\E[gg^\top]=\frac1d\E\|g\|^2I$, which holds for rotation-invariant laws. The supremum over the unit sphere is at least the average over the uniform law on it. Equality holds for $\tilde g=\Pi g$ when $k=N$. The minimum-norm solution of $\ip{E_j}{h}=y_j$ is $\Pi g$, and the ridge estimator converges to it as $\lambda\to0$. \qed

\textbf{The plain ES estimate.} For i.i.d.\ probes with $\E\hat g=g$ and $\E\|\hat g\|^2=\kappa\|g\|^2$, $\min_a\E\|a\hat g-g\|^2=(1-1/\kappa)\|g\|^2$ at $a=1/\kappa$ [N, Prop.~9]. For dense probes $\kappa=1+(d+1)/N$ (also a consequence of Theorem~\ref{thm:spike}), so the fraction of $\|g\|^2$ that is explained is $1/\kappa=N/(N+d+1)$, against the optimum $N/d$.

\subsection{Proof of Theorem~\ref{thm:spike}}
Let $z_j=\ip{E_j}{u}$ and $E_j^\perp=E_j-z_ju$. Since $E_j$ is a standard Gaussian vector, $z_j\sim N(0,1)$, $E_j^\perp$ is Gaussian with covariance $I-uu^\top$, and the two are independent. With $y_j=\|g\|z_j$,
\[
\hat g=\frac{\|g\|}{N}\sum_jz_j\big(z_ju+E_j^\perp\big)=\alpha g+\frac{\|g\|}{N}\sum_jz_jE_j^\perp,\qquad\alpha=\frac1N\sum_jz_j^2 .
\]
Given $z=(z_1,\dots,z_N)$, the sum $\sum_jz_jE_j^\perp$ is Gaussian with mean zero and covariance $\|z\|^2(I-uu^\top)$. This is the law of $\|z\|(\Gamma-\ip{\Gamma}{u}u)$ for a standard Gaussian $\Gamma$. It depends on $z$ only through $\|z\|=\sqrt{N\alpha}$, so we can take $\Gamma$ independent of $z$. Then $\frac{\|g\|}{N}\sqrt{N\alpha}=\sqrt\alpha\,\|g\|/\sqrt N$. \qed

\emph{Check.} $\E\|\hat g-g\|^2=\|g\|^2\big[\E(\alpha-1)^2+\E\alpha\,\frac{d-1}{N}\big]=\frac{d+1}{N}\|g\|^2$, the known value [2].

\subsection{A bound for finite $N$}
\begin{proposition}\label{prop:finite}
Use dense probes, let $\mathrm{rank}(G_l)\le k$ and let $T_k$ be the truncation of the SVD to the top $k$ terms. Then always $\|T_k(\Gh_l)-G_l\|_F^2\le8k\|\Gh_l-G_l\|_\op^2$, and for every $t>0$, with probability at least $1-4e^{-t^2/2}$,
\[
\|\Gh_l-G_l\|_\op\le|\alpha-1|\,s_1(G_l)+\sqrt\alpha\,\frac{\|g\|}{\sqrt N}\big(\sqrt{m}+\sqrt{n}+2t\big),\qquad\E(\alpha-1)^2=\frac2N .
\]
\end{proposition}
\begin{proof}
$T_k(\Gh_l)-G_l$ has rank at most $2k$, so its Frobenius norm is at most $\sqrt{2k}$ times its operator norm. By the triangle inequality and Weyl's inequality, $\|T_k(\Gh_l)-G_l\|_\op\le\hat s_{k+1}+\|\Gh_l-G_l\|_\op\le2\|\Gh_l-G_l\|_\op$, because $s_{k+1}(G_l)=0$. For the second claim use Theorem~\ref{thm:spike}: $\Gh_l-G_l=(\alpha-1)G_l+\sqrt\alpha\,\nu(\Gamma_l-\ip{\Gamma}{u}u_l)$, where $u_l$ is the layer-$l$ block of $u$ and $\|u_l\|_\op\le1$. A Gaussian matrix satisfies $\|\Gamma_l\|_\op\le\sqrt m+\sqrt n+t$ with probability at least $1-2e^{-t^2/2}$ [16], and $|\ip{\Gamma}{u}|\le t$ with probability at least $1-2e^{-t^2/2}$.
\end{proof}
So, up to constants, the relative error of the truncated estimate is at most $8k\big(\sqrt m+\sqrt n\big)^2\|g\|^2/(N\|G_l\|_F^2)$, compared with $mn\|g\|^2/(N\|G_l\|_F^2)$ without truncation.

\subsection{Proof of Theorem~\ref{thm:risk} (reduction to [4, 5])}
By Theorem~\ref{thm:spike}, with $X=\sum_{i\le k}x_iu_iv_i^\top$ (the scaled SVD of $G$) and $W=\Gamma_l$,
\[
\frac{\Gh}{\sqrt\alpha\,\nu\sqrt n}=\sqrt\alpha\,X+\frac{W}{\sqrt n}-\frac{\ip{\Gamma}{u}}{\sqrt n}u_l .
\]
This is the spiked model $X'+W/\sqrt n$ of [4, 5] with $X'=\sqrt\alpha X-\ip{\Gamma}{u}u_l/\sqrt n$. As $N\to\infty$, $\alpha\to1$ almost surely and $|\ip{\Gamma}{u}|/\sqrt n\to0$, so $\|X'-X\|_\op\to0$ and $X'$ has rank at most $k+1$ with singular values converging to $x_1,\dots,x_k,0$. For the model $X+W/\sqrt n$, [5, Theorems 2.8--2.9] give, for $x_i>\beta^{1/4}$,
\[
\hat s_i\to\sqrt{\Big(x_i+\frac1{x_i}\Big)\Big(x_i+\frac\beta{x_i}\Big)},\qquad
\ip{\hat u_i}{u_i}^2\to c_i^2=\frac{x_i^4-\beta}{x_i^4+\beta x_i^2},\qquad
\ip{\hat v_i}{v_i}^2\to\tilde c_i^2=\frac{x_i^4-\beta}{x_i^4+x_i^2},
\]
and for $x_i\le\beta^{1/4}$, $\hat s_i\to1+\sqrt\beta$ with asymptotically no correlation between the singular vectors. [4, Section IV] shows that the loss of a shrinker with values $\eta_i$ then converges to $\sum_i(x_i^2+\eta_i^2-2x_i\eta_ic_i\tilde c_i)$, that the minimizer is $\eta_i=x_ic_i\tilde c_i$, and that this equals the function $\eta(\hat s_i)$ of \eqref{eq:shr}. The minimum is $x_i^2(1-c_i^2\tilde c_i^2)=R(x_i)$, and a non-detected spike contributes $x_i^2$. Since these limits are continuous in the spike strengths, they are not changed by replacing $X$ with $X'$, nor by the factor $\sqrt\alpha\to1$ in the scaling. The condition $x_i>\beta^{1/4}$ reads $s_i^2N/(\|g\|^2n)>\sqrt{m/n}$, which is \eqref{eq:Nstar}. For $x\to\infty$, $(x^4-\beta)^2/(x^2(x^2+\beta)(x^2+1))=x^2-(1+\beta)+O(x^{-2})$. \qed

We state this as a reduction: the only new step is Theorem~\ref{thm:spike}, which puts the ES estimate into the model of [4, 5]. We did not re-derive the limits of [4, 5]; Figure~\ref{fig:risk}(a) checks the final formula numerically.

\subsection{Proof of Proposition~\ref{prop:cov}}
One probe gives $\Gh_1=yE$ with $E=AB^\top/\sqrt r$ and $y=\ip{E}{G}=\ip{M}{B}/\sqrt r$, $M=G^\top A\in\R^{n\times r}$. Since $\E\Gh_1=G$ and the $N$ terms are independent, $N\,\E[(\Gh-G)(\Gh-G)^\top]=\E[\Gh_1\Gh_1^\top]-GG^\top$. Now $\Gh_1\Gh_1^\top=\frac1{r^2}\ip{M}{B}^2AB^\top BA^\top$. Given $A$, Isserlis' formula for the Gaussian matrix $B$ gives
\[
\E_B\big[\ip{M}{B}^2B^\top B\big]=n\|M\|_F^2I_r+2M^\top M .
\]
With $S=GG^\top$, $\|M\|_F^2=\tr(A^\top SA)$ and $M^\top M=A^\top SA$, so $\E[\Gh_1\Gh_1^\top]=\frac1{r^2}\E\big[n\tr(A^\top SA)AA^\top+2AA^\top SAA^\top\big]$. Let $a_1,\dots,a_r$ be the columns of $A$. Using $\E[(a^\top Sa)aa^\top]=\tr(S)I+2S$ for one column and independence for two different columns,
\[
\E\big[\tr(A^\top SA)AA^\top\big]=r^2\tr(S)I+2rS,\qquad\E\big[AA^\top SAA^\top\big]=r(r+1)S+r\tr(S)I .
\]
Together, $\E[\Gh_1\Gh_1^\top]=(n+\frac2r)\tr(S)I+(\frac{2n}r+2+\frac2r)S$; subtract $GG^\top$ and note $\tr(S)=\|G\|_F^2$. The other side follows by exchanging the roles of $A$ and $B$. For dense probes, $\E[\ip{E}{G}^2EE^\top]=n\|G\|_F^2I+2GG^\top$ by Isserlis' formula. \qed

We checked the formula by Monte Carlo ($6\times9$, rank-2 $G$, $4\cdot10^5$ samples): the largest entrywise deviation, relative to the largest entry, is below $1\%$ for $r=1,2,8$.

\subsection{Proof of Proposition~\ref{prop:ceiling}}
For a quadratic, $f(W-\eta D)-f(W)=-\eta\ip{g}{D}+\frac{\eta^2}2\ip{D}{HD}$. For $D=\hat g$ with dense probes, $\E\hat g=g$ and
\[
\E\ip{\hat g}{H\hat g}=\Big(1-\frac1N\Big)\ip{g}{Hg}+\frac1N\E\big[\ip{g}{E}^2\ip{E}{HE}\big]=\Big(1+\frac1N\Big)\ip{g}{Hg}+\frac1N\tr(H)\|g\|^2,
\]
by independence of the probes and Isserlis' formula $\E[\ip{g}{E}^2\ip{E}{HE}]=\tr(H)\|g\|^2+2\ip{g}{Hg}$. Minimizing $-\eta\|g\|^2+\frac{\eta^2}2Q$ over $\eta$ gives the decrease $\|g\|^4/(2Q)$; $Q=\E\ip{\hat g}{H\hat g}$ for ES and $Q=\ip{g}{Hg}$ for the exact gradient. \qed

The second-moment identity is from our notes [N, Prop.~10]. For rank-one probes and Kronecker $H$ the term $\tr(H)\|g\|^2$ gets small corrections given there.

\section{Use of AI}
We used Claude (Anthropic) to explore ideas, to derive and check proofs, to write the code and to draft the text. The question, the earlier draft [D] and notes [N], and the final checking are ours.

\section{Experiments: details and further results}\label{app:exp}
All code is in the repository (\texttt{code/}); every number in this report is produced by the scripts named below, with fixed seeds. Only \texttt{numpy} is needed (and \texttt{scikit-learn} for the data).

\textbf{Figure~\ref{fig:risk} (\texttt{exp\_risk.py}).} $G=U\,\mathrm{diag}(s)V^\top$ with Haar factors and $s\propto(1,0.6,0.35)$. ``Shrunk'' is \eqref{eq:shr} with $\hat\nu^2=\frac1{N^2}\sum_jy_j^2$; for rank-1 and rank-4 probes with the sign-flip factor $\hat c$. Table~\ref{tab:flip} shows that the calibration is necessary for low-rank probes and costs little for dense ones.

\begin{table}[H]
\centering\small
\caption{Relative error of the shrunk estimate without and with sign-flip calibration (setting of Figure~\ref{fig:risk}). Values above 1 are worse than the zero estimate.}\label{tab:flip}
\begin{tabular}{lcccccc}
\toprule
$N$ & 60 & 160 & 400 & 1000 & 2500 & 10000\\
\midrule
dense probes, no calibration & 1.00 & 0.99 & 0.75 & 0.42 & 0.21 & 0.058\\
dense probes, calibrated & 1.00 & 1.00 & 0.83 & 0.47 & 0.23 & 0.059\\
rank-4 probes, no calibration & 21.7 & 3.41 & 1.05 & 0.52 & 0.25 & 0.068\\
rank-4 probes, calibrated & 1.00 & 1.00 & 0.86 & 0.54 & 0.27 & 0.068\\
rank-1 probes, no calibration & 70.9 & 15.1 & 4.16 & 1.14 & 0.39 & 0.097\\
rank-1 probes, calibrated & 1.00 & 1.00 & 0.99 & 0.73 & 0.37 & 0.097\\
\midrule
Theorem~\ref{thm:risk} (dense) & 1.00 & 0.99 & 0.73 & 0.42 & 0.21 & 0.058\\
plain ES, $(d+1)/N$ & 160 & 60 & 24 & 9.6 & 3.8 & 0.96\\
\bottomrule
\end{tabular}
\end{table}

\textbf{Figure~\ref{fig:descent} (\texttt{exp\_descent.py}).} $f(W)=\frac12\ip{W-W^\star}{L(W-W^\star)R}$ with diagonal $L,R$; $W^\star$ has singular values $(3,2,1)$. Stepsize grids: powers of two around $1/\kappa$ for ES and around $1$ for the shrinkage methods; the best final loss is shown. Steps to reach $10^{-2}$ of the initial loss for $H=I$: exact gradient 1, ES 50, two-level 39, shrinkage only 37 (but it stalls at $4.5\cdot10^{-3}$). Steps to $10^{-4}$: ES 99, two-level 89. Table~\ref{tab:stall} shows the stall.

\begin{table}[H]
\centering\small
\caption{Shrinkage only, stepsize 1, $H=I$, rank-one probes, one run. ``Top share'' is $s_1^2/\|G\|_F^2$ of the current gradient $G=W_t-W^\star$. ``Error'' is the relative error of the shrunk estimate at this step.}\label{tab:stall}
\begin{tabular}{lccccccccc}
\toprule
step & 0 & 1 & 2 & 3 & 4 & 6 & 8 & 10\\
\midrule
loss & 7.00 & 2.99 & 1.77 & 1.52 & 1.18 & 0.95 & 0.86 & 0.77\\
top share & 0.64 & 0.27 & 0.22 & 0.21 & 0.19 & 0.15 & 0.12 & 0.12\\
directions kept & 2 & 4 & 2 & 3 & 2 & 1 & 1 & 1\\
error & 0.43 & 0.59 & 0.86 & 0.78 & 0.89 & 0.96 & 0.98 & 0.97\\
\bottomrule
\end{tabular}
\end{table}

\textbf{Figure~\ref{fig:nn} (\texttt{exp\_nn\_wide.py}).} Network $64$--$512$--$512$--$10$, tanh, Gaussian initialization, cross-entropy on a fixed batch of 256 standardized \texttt{digits} images. The directional derivatives are exact, $y_j=\sum_la_{j,l}^\top G_lb_{j,l}$ with $G_l$ from backpropagation; this is the limit $\sigma\to0$ of an antithetic pair. In the small network, antithetic finite differences with $\sigma=10^{-3}$ differ from the exact values by $0.015\%$ in relative $\ell_2$ norm, by $1.3\%$ for $\sigma=10^{-2}$ and by $60\%$ for $\sigma=0.1$ (\texttt{check\_fd.py}); so the model \eqref{eq:model} needs a small radius.

\begin{table}[H]
\centering\small
\caption{Network $64$--$512$--$512$--$10$ after 100 gradient steps (loss $0.012$). Relative error per layer: plain ES / shrunk (measured, rank-one probes, calibrated, one draw) / Theorem~\ref{thm:risk} (dense probes).}\label{tab:nn}
\begin{tabular}{lccccc}
\toprule
layer (shape) & $N_1^\star$ & $N=1024$ & $N=4096$ & $N=16384$ & $N=65536$\\
\midrule
1 ($512\times64$) & 6641 & 257 / 1.00 / 1.00 & 65 / 1.00 / 1.00 & 16 / 0.95 / 0.92 & 3.9 / 0.64 / 0.57\\
2 ($512\times512$) & 5554 & 636 / 1.00 / 1.00 & 164 / 1.00 / 1.00 & 40 / 0.91 / 0.84 & 10 / 0.47 / 0.41\\
3 ($10\times512$) & 630 & 12 / 1.00 / 0.98 & 2.9 / 0.91 / 0.73 & 0.75 / 0.43 / 0.37 & 0.19 / 0.14 / 0.13\\
\bottomrule
\end{tabular}
\end{table}
The per-layer errors stay large in the two wide layers even when the cosine of the whole gradient is high: the estimate keeps the top directions and drops the rest, and the last layer carries about half of $\|g\|^2$.

\textbf{Small network (\texttt{exp\_nn.py}).} $64$--$96$--$96$--$10$, $d=16320$, $N=1024$ antithetic pairs with $\sigma=10^{-3}$ (true finite differences). After 300 gradient steps, $N_1^\star=2147$, $2216$, $366$. Shrinkage returns zero for layers 1 and 2 and has relative error $0.92$ in layer 3 (Theorem~\ref{thm:risk}: $1.00$, $1.00$, $0.81$). The cosine of the whole gradient is $0.24$ for ES and $0.29$ for two-level shrinkage.

TRAINTABLEPLACEHOLDER

\textbf{Iterated shrinkage on the same evaluations (\texttt{shrink.py}, \texttt{iterated\_shrink}).} We iterate $G^{(k+1)}=\mathcal S\big(G^{(k)}+\frac1N\sum_j(y_j-\ip{E_j}{G^{(k)}})E_j\big)$ with the noise level re-estimated from the residuals. This is singular value projection [7] with a data-driven threshold. For a $200\times200$ rank-3 gradient ($k(m+n)=1200$), one run each:

\begin{table}[H]
\centering\small
\caption{Relative error; ``iterated'' after 5 / 20 / 60 passes.}
\begin{tabular}{lcccc}
\toprule
 & $N=1000$ & $N=2000$ & $N=4000$ & $N=8000$\\
\midrule
rank-1 probes, one pass & 0.92 & 0.73 & 0.45 & 0.24\\
rank-1 probes, iterated & 0.89 / 0.86 / 0.87 & 0.73 / 1.00 / 1.00 & 0.45 / 1.00 / 1.00 & 0.06 / 0.09 / 0.09\\
rank-8 probes, one pass & 0.80 & 0.47 & 0.29 & 0.15\\
rank-8 probes, iterated & 0.67 / 0.47 / 0.78 & 0.50 / 0.75 / 0.82 & 0.08 / 0.66 / 0.46 & 0.003 / 0.000 / 0.000\\
\bottomrule
\end{tabular}
\end{table}
With rank-8 probes and $N=8000$ the gradient is recovered exactly from the same evaluations, as in [3]. Closer to $k(m+n)$ the plain iteration is unstable. We report this as an open point, not as a method.
